% Bagian: sets-functions-relations
% Bab: induction
% Subbagian: induction

\documentclass[../../../include/open-logic-section]{subfiles}

\begin{document}

\olfileid[id]{sfr}{ind}{int}

\olsection{Pendahuluan}

\begin{explain}
Induksi merupakan teknik pembuktian yang umum digunakan dalam matematika dan
logika. Anda mungkin telah mengenalnya dari matematika, ketika induksi
dilakukan pada bilangan asli. Ternyata, ada banyak jenis himpunan objek
matematis dan logis lain yang juga cocok untuk induksi.

Secara umum, induksi memungkinkan kita membuktikan suatu klaim universal,
yakni menunjukkan bahwa setiap objek dari jenis tertentu memiliki suatu
sifat. Secara khusus, induksi sering berguna apabila metode bukti
``introduksi universal'' terlalu rumit. Induksi hanya berlaku pada objek
matematis yang dibangun dengan cara khusus: jika setiap unsur dalam himpunan
bersifat dasar atau dibangun dari unsur-unsur dasar, maka kita dapat
menggunakan induksi pada himpunan itu. Secara lebih tepat:
\end{explain}

\begin{defn}
Suatu himpunan $S$ disebut \emph{tertutup} terhadap fungsi $f$ jika dan hanya
jika, setiap kali $x \in S$, berlaku $f(x) \in S$.
\end{defn}

\begin{explain}
Sebagai contoh, bilangan asli ($\mathbb{N}$) tertutup terhadap fungsi penerus
$f(x) = x+1$. Jika $x$ bilangan asli, maka $x+1$ juga bilangan asli. Namun,
bilangan asli tidak tertutup terhadap fungsi pendahulu $f(x) = x-1$, karena
$0$ bilangan asli, sedangkan $-1$ bukan.
\end{explain}

\begin{defn}
Suatu sifat $P$ dikatakan \emph{dipertahankan oleh} fungsi $f(x)$ apabila,
untuk setiap objek $a$ dalam domain $f$, $P(a)$ mengimplikasikan $P(f(a))$
(jika $a$ memiliki sifat $P$, maka $f(a)$ juga demikian).
\end{defn}

\begin{explain}
Sifat ``merupakan bilangan genap'' dipertahankan oleh fungsi $f(x) = x+2$,
karena jika $x$ genap, maka $x+2$ juga genap. Sifat merupakan bilangan genap
tidak dipertahankan oleh fungsi penerus, karena jika $x$ genap, maka $x+1$
ganjil.

Induksi berlaku pada bilangan asli karena setiap bilangan asli dapat dicapai
dari 0 dengan menerapkan fungsi penerus sebanyak berhingga kali. Hal ini
merupakan ciri setiap himpunan yang padanya kita dapat melakukan induksi.
\end{explain}

\begin{defn}[Induksi pada $\Nat$]
Jika 0 memiliki suatu sifat $P$, dan jika $P$ dipertahankan oleh fungsi
penerus, maka setiap bilangan asli memiliki sifat $P$.
\end{defn}

\begin{ex} Jumlah bilangan asli dari 0 sampai $n$ ialah $\frac{n(n+1)}{2}$. \\

\textbf{Kasus dasar.} Untuk $n=0$, jumlah tersebut ialah $0=
\frac{0\cdot 1}{2}$.\\

\textbf{Langkah induktif.} Misalkan sifat tersebut berlaku bagi $k$. Kita
akan menunjukkan bahwa sifat itu berlaku bagi $k+1$. Dengan kata lain,
misalkan $P(k)$ berlaku, yakni,
\[ 0 + 1 + \ldots + k = \frac{k(k+1)}{2} \]
Kita akan menunjukkan bahwa $P(k+1)$ berlaku, yakni bahwa
\[ 0 + 1 + \ldots + k + (k+1) = \frac{(k+1)(k+2)}{2} \]
dengan menggunakan asumsi bahwa $P(k)$ berlaku.

\begin{align*}
0 + 1 + \ldots + k &= \frac{k(k+1)}{2} \tag{Asumsi} \\
0 + 1 + \ldots + k + (k+1) &= \frac{k(k+1)}{2} + (k+1)
\tag{Tambahkan $k+1$ pada kedua ruas} \\
&= \frac{k(k+1) + 2(k+1)}{2} \\
&= \frac{k^2 + 3k + 2}{2} \\
&=\frac{(k+1)(k+2)}{2}
\end{align*}
Jadi, langkah induktif berhasil, dan klaim tersebut telah kita buktikan.
\end{ex}

\begin{explain}
Untuk !!{formula}s, unsur-unsur dasarnya ialah !!{formula}s atomik.
!!{formula}s yang lebih kompleks diperoleh dengan menggabungkan
!!{formula}s yang lebih sederhana menggunakan operator. Setiap operator dapat
dipandang sebagai fungsi pada !!{formula}s yang menghasilkan !!{formula}
baru dengan menggabungkan kedua formula tersebut menggunakan operator yang
bersangkutan. Sebagai contoh, fungsi yang menggabungkan dua !!{formula}s
$!!^a$ dan $!B$ menjadi suatu konjungsi dapat ditulis sebagai
$\mathcal E _\land (!!^a, !B) = !!^a \land !B$. Kita dapat menyebut fungsi
semacam itu \emph{fungsi pembentuk formula}. Himpunan !!{formula}s tertutup
terhadap fungsi-fungsi ini: setiap kali $!!^a$ dan $!B$ merupakan
!!{formula}s, demikian pula $\lnot !!^a$, $!!^a \land !B$, dan sebagainya.
\end{explain}

\begin{defn}[Induksi pada !!{formula}s]
Jika setiap !!{formula} atomik memiliki sifat $P$ dan $P$ dipertahankan oleh
fungsi-fungsi pembentuk !!{formula}, maka setiap !!{formula} memiliki sifat
$P$.
\end{defn}

\begin{explain}
Definisi ini menyarankan suatu ``resep'' bagi bukti induktif pada
!!{formula}s. Jika kita diminta menunjukkan bahwa setiap !!{formula} memiliki
sifat $P$, ikuti langkah-langkah berikut:\\

\textbf{Kasus Dasar.} Misalkan $!!^a$ suatu !!{formula} atomik. [\ldots]
Jadi, $!!^a$ memiliki sifat $P$.

\textbf{Langkah Induktif.} Misalkan $!!^a$ dan $!B$ merupakan !!{formula}s
yang keduanya memiliki sifat $P$.

Kasus $\lnot$: [\ldots] Jadi, $\lnot !!^a$ memiliki sifat $P$.

Kasus $\land$: [\ldots] Jadi, $!!^a \land !B$ memiliki sifat $P$.

Kasus $\lor$: [\ldots] Jadi, $!!^a \lor !B$ memiliki sifat $P$.

Kasus $\lif$: [\ldots] Jadi, $!!^a \lif !B$ memiliki sifat $P$.

Kasus $\liff$: [\ldots] Jadi, $!!^a \liff !B$ memiliki sifat $P$.

Kasus $\lforall[]$: [\ldots] Jadi, $\lforall[x] !!^a$ memiliki sifat $P$.

Kasus $\lexists[]$: [\ldots] Jadi, $\lexists[x] !!^a$ memiliki sifat $P$. \\

Jadi, setiap !!{formula} memiliki sifat $P$.
\end{explain}

\end{document}
